\chapter{Homeopathy}
\index{Homeopathy}

\section*{Definition and Overview}
Homeopathy is a system of complementary medicine based on two principles: "like cures like", meaning a substance that causes symptoms in a healthy person can treat similar symptoms in a sick person, and the law of infinitesimals, which holds that the more a remedy is diluted, the more potent it becomes. Homeopathic remedies are prepared by repeatedly diluting plant, mineral, or animal substances, often to the point where no molecules of the original substance remain. Extensive scientific evidence shows that homeopathic remedies work no better than placebo for most conditions, and major health authorities do not recommend it as a treatment for serious illness.

\section*{Epidemiology}
Homeopathy is used by a small but significant minority of people in many countries, often alongside conventional medicine. Around 2--5\% of people have used homeopathic remedies, with higher use among those interested in natural therapies and among parents treating children. people spend millions of dollars annually on homeopathic products. Use is higher in Europe, where homeopathy has a long tradition. Most use is for minor, self-limiting conditions such as colds, allergies, and anxiety, and most people also use conventional care.

\section*{Aetiology and Pathophysiology}
\begin{itemize}
    \item \textbf{The principle of similars:} substances that cause symptoms are used to treat similar symptoms
    \item \textbf{Potentisation:} serial dilution and vigorous shaking (succussion) are claimed to increase potency
    \item \textbf{Ultra-dilution:} most remedies are diluted beyond Avogadro's number, meaning no molecules of the original substance remain
    \item \textbf{The placebo effect:} the reported benefits of homeopathy are attributed by scientists to placebo, expectation, natural recovery, and regression to the mean
    \item \textbf{No plausible mechanism:} there is no scientifically accepted mechanism by which such dilutions could act
    \item \textbf{Regulation:} in many countries, homeopathic products are regulated as medicines by the TGA, but efficacy is not evaluated for most products
    \item \textbf{Evidence reviews:} systematic reviews find homeopathy no better than placebo for any condition
\end{itemize}

\section*{Clinical Features}
\begin{itemize}
    \item People may seek homeopathy for chronic conditions, anxiety, allergies, and childhood illnesses
    \item Remedies are available as tablets, granules, liquids, and creams
    \item "Individualisation": practitioners choose remedies based on a detailed consultation about physical and emotional symptoms
    \item Reported improvements are common and may reflect natural recovery or placebo
    \item Risks are low when remedies are used alone for minor conditions
    \item The main risk is replacing effective treatment with ineffective remedies
\end{itemize}

\section*{Differential Diagnosis}
\begin{itemize}
    \item The apparent benefit of homeopathy must be distinguished from the natural course of the illness
    \item Symptoms attributed to homeopathy may reflect the underlying condition
    \item Delayed or abandoned conventional treatment is the principal danger
    \item Some herbal or homeopathic products contain active ingredients that can cause harm
    \item Serious illness must be treated with evidence-based care
\end{itemize}

\section*{When to Seek Medical Care}
Anyone with a serious, persistent, or worsening condition should seek conventional medical care, and no one should delay or stop effective treatment in favour of homeopathy. Tell your doctor about all complementary products you use.

\textbf{Contact your local emergency services for:}
\begin{itemize}
    \item Severe illness in a child, including high fever, breathing difficulty, or drowsiness
    \item Chest pain, severe abdominal pain, or collapse
    \item Symptoms of serious infection: stiff neck, severe headache, or rash
    \item Any condition that is worsening despite homeopathic treatment
    \item Allergic or toxic reactions to a product
\end{itemize}

\section*{Diagnosis and Investigations}
\begin{itemize}
    \item \textbf{Assessment of the underlying condition:} conventional diagnosis is essential
    \item \textbf{Medicine review:} all complementary products are documented
    \item \textbf{Monitoring:} objective measures of progress, such as blood pressure, blood tests, or imaging, where relevant
    \item \textbf{Evidence appraisal:} review of what is known about the product and condition
    \item \textbf{Risk assessment:} whether delaying effective treatment poses harm
\end{itemize}

\section*{Management}
\begin{itemize}
    \item \textbf{Evidence-based treatment:} conventional care remains the standard for all conditions
    \item \textbf{Discussion:} health professionals respectfully explore the reasons for using homeopathy
    \item \textbf{Safe integration:} if the patient chooses to use homeopathy, it should not replace effective treatment
    \item \textbf{Monitoring:} the condition is followed objectively
    \item \textbf{Education:} clear information about the evidence helps informed decisions
    \item \textbf{Regulation awareness:} consumers are advised to choose products from reputable sources
\end{itemize}

\section*{Complications}
\begin{itemize}
    \item Delayed diagnosis and treatment of serious conditions
    \item Discontinuation of effective medicine, leading to disease progression
    \item Serious outcomes in children and in conditions such as infections, diabetes, and cancer
    \item Financial cost of ineffective products
    \item Rare toxic effects from poorly prepared or contaminated products
    \item Misplaced confidence leading to missed opportunities for cure
\end{itemize}

\section*{Prognosis and Outcomes}
Homeopathic remedies themselves carry no direct physiological benefit beyond placebo. Outcomes for users are therefore determined by the underlying condition: minor, self-limiting illnesses improve regardless of treatment, while serious conditions require effective medical care. People who use homeopathy alongside conventional treatment for minor complaints generally come to no harm, whereas those who substitute it for proven treatment can have poor outcomes. Evidence-based medicine offers the most reliable path to recovery.

\section*{Prevention and Self-Care}
\begin{itemize}
    \item Use evidence-based treatments for all health conditions
    \item Never stop prescribed medicines without medical advice
    \item Seek conventional care for serious, persistent, or worsening symptoms
    \item Ask your doctor before using any complementary product
    \item Be wary of remedies promising cures for serious disease
    \item Keep children with significant illness under medical care
    \item Save money on products without proven benefit
    \item Discuss any complementary therapy use openly with your health team
\end{itemize}

\section*{Sources and Further Reading}
The following reputable sources provide further information for healthcare students and patients seeking reliable, current advice about homeopathy.

\begin{itemize}
    \item Therapeutic Goods Administration. \textit{Complementary medicines regulation.} https://www.tga.gov.au/
    \item National Centre for Complementary and Integrative Health. \textit{Homeopathy: what you need to know.} https://www.nccih.nih.gov/
    \item Australian Skeptics and Friends of Science in Medicine. \textit{Evidence appraisal of homeopathy.} https://www.skeptics.com.au/
    \item NHS. \textit{Homeopathy: overview and evidence.} https://www.nhs.uk/conditions/homeopathy/
    \item Cochrane Library. \textit{Systematic reviews of homeopathy.} https://www.cochranelibrary.com/
    \item Australian Government National Health and Medical Research Council. \textit{NHMRC homeopathy review.} https://www.nhmrc.gov.au/
\end{itemize}
